% !TEX root = ../tex/7-4.tex
% 第7章：本文は青、補足は緑
\newcommand{\asmain}{%
    \setbeamercolor{frametitle}{fg=white,bg=main}%
    \setbeamercolor{structure}{fg=main}%
    \setbeamercolor{itemize item}{fg=main}%
    \setbeamercolor{itemize subitem}{fg=main}%
    \setbeamercolor{enumerate item}{fg=main}%
}
\newcommand{\assuppl}{%
    \setbeamercolor{frametitle}{fg=white,bg=examplecol}%
    \setbeamercolor{structure}{fg=examplecol}%
    \setbeamercolor{itemize item}{fg=examplecol}%
    \setbeamercolor{itemize subitem}{fg=examplecol}%
    \setbeamercolor{enumerate item}{fg=examplecol}%
}
\newcommand{\supplmark}{%
    \setlength{\fboxsep}{1.5pt}%
    \setlength{\fboxrule}{0.7pt}%
    \fbox{\textbf{補足}}%
}
\asmain


%=============================================
\asmain
\begin{frame}{7.4 の問い}
    \begin{block}{節タイトル}
        Is Diffusion Model's SDE Optimal Solution to SB Problem?
    \end{block}
    \begin{itemize}
        \item 短答：\textbf{一般には No}（両端点を厳密に結ぶフル SB ではない）
        \item ただし：ポテンシャル $\Psi\equiv 1$ の極限では SB の特別ケースに見える
        \item より正確には：\textbf{Schrödinger half-bridge}（片側制約）として読む
    \end{itemize}
\end{frame}
%=============================================
\assuppl
\begin{frame}{\supplmark：7.4 の位置づけ（Optional 章）}
    \begin{block}{この節の役割（論文 7.4）}
        7.3 の関係図の上に、標準的なスコアベース拡散を置く。\\
        「拡散 $=$ フル SB の最適解か？」を Yes/No で切り分ける。
    \end{block}
    \begin{enumerate}
        \item \textbf{7.4.1} 拡散は SB の特別ケース（$\Psi\equiv 1$）
        \item \textbf{7.4.2} 半橋（half-bridge）と端点制約の差；DSB/IPF
    \end{enumerate}
\end{frame}
%=============================================
\asmain
\begin{frame}{7.4 の道筋}
    \begin{enumerate}
        \item SB のポテンシャル分解と制御ドリフト（再掲）
        \item $\Psi\equiv 1$ で Anderson 逆時間 SDE が出る
        \item 境界条件が合うとき／合わないとき
        \item half-bridge：片側だけ $P_0$ または $P_T$ を合わせる
        \item 拡散：$p_T\approx p_{\mathrm{prior}}$（$T\to\infty$）vs SB：有限 $T$ で厳密一致
        \item DSB：IPF でフル橋へ近づける
    \end{enumerate}
\end{frame}
%=============================================
\section{7.4.1 拡散は SB の特別ケース}
\asmain
\begin{frame}{区切り：7.4.1 Diffusion as a Special Case of SB}
    \tableofcontents[currentsection]
\end{frame}
%=============================================
\asmain
\begin{frame}{SB が拡散より広い理由}
    SB は任意の $p_{\mathrm{src}},p_{\mathrm{tgt}}$ のあいだを、\\
    参照拡散に\textbf{制御ドリフト}を足して橋渡しする。
    \begin{block}{ポテンシャル分解（再掲）}
        \[
            \nabla\log\Psi(x,t)+\nabla\log\widehat\Psi(x,t)
            =
            \nabla\log p_t^{\mathrm{SB}}(x).
        \]
        $\Psi,\widehat\Psi$ が式 (7.2.12) の Schrödinger 系を満たす。
    \end{block}
    非自明なポテンシャル $\Rightarrow$ ガウス事前以外の端点にも曲がれる。
\end{frame}
%=============================================
\asmain
\begin{frame}{$\Psi\equiv 1$ の還元}
    ポテンシャルを定数 $\Psi(x,t)\equiv 1$ とすると、\\
    (7.2.12) の第二式は通常の Fokker--Planck になり
    \[
        \widehat\Psi(x,t)=p_t^{\mathrm{SB}}(x).
        \tag{7.4.1}
    \]
    \begin{itemize}
        \item 前向き SB-SDE $=$ \textbf{無制御}の参照過程
        \[
            dx_t=f(x_t,t)\,dt+g(t)\,dw_t
        \]
        \item 逆向きはスコア補正のみ（次スライド）
    \end{itemize}
\end{frame}
%=============================================
\asmain
\begin{frame}{Anderson の逆時間 SDE との一致}
    $\Psi\equiv 1$ のとき逆向き SB-SDE は
    \[
        dx_t
        =
        \bigl[
            f(x_t,t)
            -
            g^{2}(t)\nabla\log p_t^{\mathrm{SB}}(x_t)
        \bigr]dt
        +
        g(t)\,d\bar w_t.
    \]
    \begin{exampleblock}{対応}
        これは拡散モデルで使う Anderson の reverse-time SDE と同じ形。\\
        つまり拡散は、\textbf{追加制御なし}（ゼロ制御極限）の SB に見える。
    \end{exampleblock}
\end{frame}
%=============================================
\asmain
\begin{frame}{ただし境界条件が鍵}
    上の還元は、式 (7.2.12) の境界
    \[
        \Psi(\cdot,0)\widehat\Psi(\cdot,0)=p_{\mathrm{src}},
        \qquad
        \Psi(\cdot,T)\widehat\Psi(\cdot,T)=p_{\mathrm{tgt}}
    \]
    が $\Psi\equiv 1$ と両立するときに限って「フル SB」になる。
    \begin{alertblock}{一般の $(p_{\mathrm{src}},p_{\mathrm{tgt}})$}
        $\Psi\equiv 1$ では境界を満たせない。\\
        フル SB は非自明な $\Psi,\widehat\Psi$ で参照を曲げ、両端を厳密に合わせる。
    \end{alertblock}
\end{frame}
%=============================================
\asmain
\begin{frame}{拡散モデル側の設計}
    \begin{columns}[T,onlytextwidth]
        \begin{column}{0.48\textwidth}
            \begin{block}{拡散}
                一端を簡単な事前\\
                （典型的にはガウス）に固定。\\
                学習するのは逆時間スコア。
            \end{block}
        \end{column}
        \begin{column}{0.48\textwidth}
            \begin{block}{フル SB}
                両端を任意に指定。\\
                ポテンシャル由来の\\
                非線形制御ドリフトが必要。
            \end{block}
        \end{column}
    \end{columns}
    \vfill
    傘としては SB の方が広い：非自明ポテンシャルで任意端点、\\
    $\Psi\equiv 1$ で拡散モデルケースへ潰れる。
\end{frame}
%=============================================
\asmain
\begin{frame}{線形拡散での注意}
    標準の線形（VP など）拡散では
    \[
        p_T\approx p_{\mathrm{prior}}
        \quad\text{は}\quad
        T\to\infty
    \]
    の漸近でしか成り立たないことが多い。
    \begin{itemize}
        \item 有限 $T$ では $p_T$ はなお $p_{\mathrm{data}}$ に依存しうる
        \item したがって「事前への厳密マッチ」は SB ほど強くない
    \end{itemize}
\end{frame}
%=============================================
\section{7.4.2 Schrödinger half-bridges}
\asmain
\begin{frame}{区切り：7.4.2 Diffusion as Schrödinger Half-Bridges}
    \tableofcontents[currentsection]
\end{frame}
%=============================================
\asmain
\begin{frame}{なぜ「フル橋」ではないか}
    \begin{block}{結論の先取り}
        拡散モデルは両端制約の SB 最適解というより、\\
        \textbf{片側だけ合わせる half-bridge} として理解するのが正確。
    \end{block}
    フル SB：参照 $R$ に近く、かつ $P_0=p_{\mathrm{src}}$ かつ $P_T=p_{\mathrm{tgt}}$。\\
    半橋：どちらか一方だけを課す緩和。
\end{frame}
%=============================================
\asmain
\begin{frame}{half-bridge の定義}
    参照パス測度を $R$ とする。
    \begin{itemize}
        \item \textbf{前向き半橋}：
        \[
            \min_{P:\,P_0=p_{\mathrm{src}}}
            D_{\mathrm{KL}}(P\|R)
        \]
        \item \textbf{後ろ向き半橋}：
        \[
            \min_{P:\,P_T=p_{\mathrm{tgt}}}
            D_{\mathrm{KL}}(P\|R)
        \]
    \end{itemize}
    \begin{exampleblock}{言葉}
        「望む初期（または終端）から出発（到着）する過程のうち、\\
        参照に一番似ているものは？」
    \end{exampleblock}
\end{frame}
%=============================================
\asmain
\begin{frame}{反復でフル橋へ}
    前向き半橋と後ろ向き半橋を交互に解くと、\\
    Iterative Proportional Fitting (IPF) 的にフル SB を近似できる。
    \begin{alertblock}{計算の難しさ}
        両境界を同時に厳密に課すのは重い。\\
        半橋はその緩和であり、拡散の学習手順とも相性が良い。
    \end{alertblock}
\end{frame}
%=============================================
\asmain
\begin{frame}{端点マッチの差}
    標準拡散の前向き SDE は線形で、$T\to\infty$ で
    \[
        p_T\approx p_{\mathrm{prior}}.
    \]
    有限時刻では $x_T=\alpha_T x_0+\sigma_T\epsilon$ の周辺なので、\\
    一般に $p_T$ は $p_{\mathrm{data}}$ 依存のまま（付録 C.1.5）。
    \vfill
    SB は制御 $g^{2}(t)\nabla_x\log\Psi$ により、有限 $T$ で
    \[
        p_T=p_{\mathrm{prior}}
    \]
    を\textbf{厳密に}課せる。
\end{frame}
%=============================================
\asmain
\begin{frame}{対比表}
    \renewcommand{\arraystretch}{1.35}
    \begin{tabular}{lll}
        & \textbf{拡散モデル} & \textbf{Schrödinger Bridge} \\
        \hline
        終端周辺
        & $p_T\approx p_{\mathrm{prior}}$（$T\to\infty$）
        & $p_T=p_{\mathrm{prior}}$（有限 $T$）
        \\
        追加制御
        & 基本なし（参照＋スコア逆過程）
        & $\Psi,\widehat\Psi$ 由来のドリフト
        \\
        制約
        & 片側に近い／漸近
        & 両端を厳密指定
        \\
    \end{tabular}
\end{frame}
%=============================================
\asmain
\begin{frame}{拡散 $=$ 前向き半橋の見方}
    固定した前向き noising 過程は、\\
    $p_{\mathrm{data}}$ にアンカーした\textbf{前向き half-bridge} と読める。
    \begin{itemize}
        \item $P_0=p_{\mathrm{data}}$ は満たす
        \item 終端は参照が生む周辺であり、一般に $p_{\mathrm{prior}}$ と厳密一致しない
        \item よって指定ペア $(p_{\mathrm{data}},p_{\mathrm{prior}})$ のフル橋ではない
    \end{itemize}
\end{frame}
%=============================================
\asmain
\begin{frame}{Diffusion Schrödinger Bridge (DSB)}
    De Bortoli et al.\ (2021)：IPF の考えで両端を交互に合わせ、\\
    拡散をフル SB ソルバへ拡張する。
    \begin{block}{Step 0}
        $P^{(0)}:=R_{\mathrm{fwd}}$（参照前向き SDE、$x_0\sim p_{\mathrm{data}}$）。\\
        $P_0^{(0)}=p_{\mathrm{data}}$ だが、通常 $P_T^{(0)}\neq p_{\mathrm{prior}}$。
    \end{block}
\end{frame}
%=============================================
\asmain
\begin{frame}{DSB：後ろ向きパス}
    \begin{block}{Step 1}
        \[
            P^{(1)}
            =
            \argmin_{P:\,P_T=p_{\mathrm{prior}}}
            D_{\mathrm{KL}}(P\|P^{(0)}).
        \]
        スコア網 $s_{\phi^\times}$ で逆時間 SDE を近似し、\\
        $x_T\sim p_{\mathrm{prior}}$ から後ろ向きシミュレート。
    \end{block}
    これで $P_T^{(1)}=p_{\mathrm{prior}}$ だが、一般に $P_0^{(1)}\neq p_{\mathrm{data}}$。
\end{frame}
%=============================================
\asmain
\begin{frame}{DSB：交互反復}
    \begin{enumerate}
        \item 前向きを学習し直し、$P_0$ を $p_{\mathrm{data}}$ へ戻す
        \item 再び後ろ向きで $P_T=p_{\mathrm{prior}}$ を課す
        \item 収束まで繰り返し、両端を満たす $P^*$ へ
    \end{enumerate}
    温和な条件で収束が示されている（De Bortoli et al., 2021）。\\
    枠組み自体は任意の端点ペアに使える。
\end{frame}
%=============================================
\assuppl
\begin{frame}{\supplmark：DSB の実務コスト}
    \begin{itemize}
        \item IPF の各段で別のスコアマッチングが必要
        \item 反復を重ねると近似誤差が蓄積しうる
        \item 安定・スケール改善や、反復依存を減らす SB 手法も多い\\
        （例：DSBM, GSBM など；本文は文献ポインタに留める）
    \end{itemize}
    \begin{exampleblock}{ゼミでの位置}
        「標準拡散は半橋；フル橋にするには追加の交互投影が要る」\\
        が本筋。実装詳細は後続文献へ。
    \end{exampleblock}
\end{frame}
%=============================================
\section{タイトルへの答え}
\asmain
\begin{frame}{タイトルへの答え（整理）}
    \begin{block}{いつ Yes に近いか}
        $\Psi\equiv 1$ で境界が両立し、参照が拡散の前向き過程と一致するとき。\\
        形式的には「ゼロ制御の SB」＝ Anderson 逆時間 SDE。
    \end{block}
    \begin{alertblock}{いつ No か（通常の実務設定）}
        両端 $(p_{\mathrm{data}},p_{\mathrm{prior}})$ を有限 $T$ で厳密に課すフル SB としては No。\\
        $p_T\approx p_{\mathrm{prior}}$ は漸近であり、制御ポテンシャルも学習しない。
    \end{alertblock}
\end{frame}
%=============================================
\asmain
\begin{frame}{答えの一文版}
    \begin{enumerate}
        \item \textbf{特別ケースとしては Yes}：\\
        SB の $\Psi\equiv 1$（無追加制御）極限が拡散 SDE／逆 SDE。
        \item \textbf{最適フル橋としては No}：\\
        標準拡散は片側／漸近制約の half-bridge。
        \item \textbf{フル SB へ行くには}：\\
        非自明ポテンシャル、または DSB のような IPF 反復が必要。
    \end{enumerate}
\end{frame}
%=============================================
\section{まとめ}
\asmain
\begin{frame}{7.4 のまとめ}
    \begin{itemize}
        \item SB は任意端点への制御付き橋；拡散はその傘の下の特殊形
        \item $\Psi\equiv 1$ $\Rightarrow$ 参照前向き＋Anderson 逆時間 SDE
        \item 境界非両立ならフル SB にはならない
        \item half-bridge：一端のみ；拡散の noising はその一例
        \item 端点：$p_T\approx p_{\mathrm{prior}}$（拡散）vs 有限 $T$ で厳密（SB）
        \item DSB/IPF：交互投影でフル橋へ
    \end{itemize}
\end{frame}